% This file is part of the modular source for Stratified Relational Workspace Theory.
% !TeX root = ../Stratified_Relational_Workspace_Theory.tex
\theorypart{XVII}{Decomposition, Coarse-Graining, and Robust Relationality}

Relationality is always assessed relative to a way of dividing a realization
into components. That dependence is not a minor technicality. If every
participant, artifact, and channel is merged into one component, the complete
workspace becomes local by definition. If a centralized state is recoded into
several complementary shares, the same centralized process can be made to look
irreducibly distributed. A defensible theory must therefore say which
partitions count as admissible, which quantities survive faithful aggregation,
and how sensitive a certificate is to alternative descriptions of the same
system.

\topic{54}{The Decomposition Problem}

Let \(\Omega\) be a finite set of primitive carriers at a declared audit
resolution. A primitive carrier may be a person, a model instance, a process,
a memory store, an instrument, or another independently tracked mechanism.
The word \emph{primitive} is relative to the audit. It does not assert a final
metaphysical atomization of the system.

\begin{definition}[System decomposition]\label{def:system-decomposition}
A decomposition of a carrier presentation is a tuple
\[
\mathcal D
=
\bigl(
\Pi,\mathbf O,G_{\Pi},\mathcal I_{\Pi}
\bigr),
\]
where:
\begin{itemize}
\item \(\Pi=\{B_\alpha\}_{\alpha\in A}\) is a partition of \(\Omega\);
\item \(O_\alpha\) is the declared local-access map for block \(B_\alpha\),
 producing the local state or history available at that component;
\item \(G_{\Pi}\) records the declared communication, transfer, and causal
 interfaces among the blocks; and
\item \(\mathcal I_{\Pi}\) is the family of feasible interventions at that
 resolution.
\end{itemize}
The accessible information of component \(\alpha\) is the sigma-field
\[
\mathscr F_\alpha^{\mathcal D}
=
\sigma\!\left(O_\alpha(\mathbf H_t)\right).
\]
Every shared bus, common memory, supervisor, or artifact that can alter more
than one local state must appear either as its own carrier or as an explicit
channel in \(G_{\Pi}\).
\end{definition}

The same physical realization can support several decompositions. One may
separate a model from its executor, merge them into one computational unit, or
treat a persistent artifact as its own component. These alternatives are not
mere renamings when they change local access or feasible interventions. Every certificate should therefore display the full specification
\[
\SRW(W_t;\varsigma),
\qquad
\varsigma=(\mathcal D,Y^+,\ell,\tau,\mathcal I,\mathcal N,\Theta).
\]
This part varies the decomposition coordinate while holding the remaining
coordinates fixed. Part~XVIII treats the full specification family.

\begin{definition}[Refinement and coarsening]\label{def:decomposition-refinement}
Let \(\mathcal D_f\) and \(\mathcal D_c\) be decompositions of the same carrier
presentation. The fine decomposition refines the coarse decomposition,
written
\[
\mathcal D_f\preceq\mathcal D_c,
\]
when there is a surjection
\[
q:V_f\longrightarrow V_c
\]
such that each coarse component is the union of the fine components in one
fiber of \(q\), and its accessible information is generated only from those
fine components. The aggregation is \emph{information faithful} when
\[
\mathscr F_B^{\mathcal D_c}
=
\bigvee_{i\in q^{-1}(B)}
\mathscr F_i^{\mathcal D_f}
\qquad(B\in V_c).
\]
It is \emph{structurally faithful} when it also retains every internal
compatibility condition and every boundary transfer condition needed to
recover the original global-section problem.
\end{definition}

Refinement and coarsening are directional. A coarse component may contain
information that was spread across several finer nodes. Conversely, a fine
description may expose internal channels that disappear when a macro-node is
formed. The next results identify what is preserved when the aggregation is
faithful.

\topic{55}{Faithful Structural Coarse-Graining}

Let the vertex set of a core-transfer diagram be partitioned as
\[
V=V_1\sqcup\cdots\sqcup V_m.
\]
Write the transfer discrepancy as
\[
D_H
=
\begin{pmatrix}
D_{\mathrm{int}}\\
D_{\partial}
\end{pmatrix},
\qquad
D_{\mathrm{int}}
=
\bigoplus_{\alpha=1}^{m}D_\alpha,
\]
where \(D_\alpha\) contains the edge constraints internal to \(V_\alpha\),
and \(D_{\partial}\) contains the constraints crossing between blocks. Define
the internally coherent macro-state spaces
\[
W_\alpha=\ker D_\alpha,
\qquad
W_{\Pi}=\bigoplus_{\alpha=1}^{m}W_\alpha.
\]
Let
\[
J_\Pi:W_\Pi\hookrightarrow\mathcal C_V
\]
be the canonical isometric inclusion, and define the quotient discrepancy
\[
D_\Pi=D_{\partial}J_\Pi.
\]
The quotient does not replace a group of nodes by an arbitrary summary. Its
macro-state is the complete internal section space of that group, and its
boundary operator retains all cross-block compatibility equations.

\begin{theorem}[Faithful coarse-graining]\label{thm:faithful-coarse-graining}
For every finite partition of the vertex set,
\[
\boxed{
\ker D_H
=
J_\Pi\ker D_\Pi.
}
\]
Equivalently,
\[
\boxed{
\Gamma(H)
\cong
\ker D_\Pi
=
\ker(D_\Pi^*D_\Pi).
}
\]
In particular,
\[
\dim\Gamma(H)
=
\dim W_\Pi-\operatorname{rank}(D_\Pi).
\]
Thus faithful coarse-graining preserves the exact structural workspace even
though it changes the number and interpretation of the displayed nodes.
\end{theorem}

\begin{proof}
A vector \(c\in\mathcal C_V\) lies in \(\ker D_H\) exactly when
\[
D_{\mathrm{int}}c=0
\qquad\text{and}\qquad
D_{\partial}c=0.
\]
The first condition is equivalent to \(c=J_\Pi w\) for a unique
\(w\in W_\Pi\). Substitution into the second condition gives
\[
D_\Pi w
=
D_{\partial}J_\Pi w
=0.
\]
Hence \(c\in\ker D_H\) if and only if
\(c=J_\Pi w\) with \(w\in\ker D_\Pi\). The Laplacian identity follows from
\(\ker(D_\Pi^*D_\Pi)=\ker D_\Pi\), and the dimension formula is rank-nullity.
\end{proof}

The theorem supplies a precise sense in which a workspace may be represented
at several resolutions without changing its exact core. The preservation
fails when the macro-state discards internal section coordinates or when
boundary equations are omitted.

\begin{corollary}[Two-scale coherence bound]\label{cor:two-scale-decomposition}
Let \(P_\Pi\) be the orthogonal projector onto \(J_\Pi W_\Pi\). Let
\(\gamma_{\mathrm{int}}>0\) be the smallest positive eigenvalue of
\(D_{\mathrm{int}}^*D_{\mathrm{int}}\), and let
\(\gamma_\Pi>0\) be the smallest positive eigenvalue of
\(D_\Pi^*D_\Pi\). Then every \(c\in\mathcal C_V\) satisfies
\[
\boxed{
\operatorname{dist}(c,\Gamma(H))
\leq
\frac{\|D_{\mathrm{int}}c\|}{\sqrt{\gamma_{\mathrm{int}}}}
+
\frac{
\|D_\Pi J_\Pi^*P_\Pi c\|
}{\sqrt{\gamma_\Pi}}.
}
\]
The first term measures incoherence within macro-components. The second
measures incompatibility among their internally coherent states.
\end{corollary}

\begin{proof}
By the triangle inequality,
\[
\operatorname{dist}(c,\Gamma(H))
\leq
\|c-P_\Pi c\|
+
\operatorname{dist}(P_\Pi c,\Gamma(H)).
\]
The workspace Poincar\'e estimate applied to \(D_{\mathrm{int}}\) bounds the
first term. Write \(P_\Pi c=J_\Pi w\). Theorem~\ref{thm:faithful-coarse-graining}
and the corresponding estimate for \(D_\Pi\) bound the second term by
\(\|D_\Pi w\|/\sqrt{\gamma_\Pi}\). Since \(J_\Pi\) is isometric,
\(w=J_\Pi^*P_\Pi c\).
\end{proof}

This decomposition is useful computationally. Internal blocks can be solved
or compressed independently, after which only the boundary discrepancy must
be assembled globally. It also gives a diagnostic distinction between a
workspace that fails inside its local components and one whose local
components are coherent but mutually incompatible.

\topic{56}{Predictive Risk Under Aggregation}

The structural core can be preserved by faithful coarse-graining, while
coalition counts generally are not. Let
\[
\mathscr F_S^{\mathcal D}
=
\bigvee_{i\in S}\mathscr F_i^{\mathcal D}
\]
be the information available to coalition \(S\) under decomposition
\(\mathcal D\). For a declared loss \(\ell\), write
\[
\Risk_{\mathcal D}(S)
=
\inf_f
\mathbb E\!\left[
\ell\bigl(Y^+,f(\mathscr F_S^{\mathcal D})\bigr)
\right].
\]

\begin{theorem}[Bayes-risk transport under faithful aggregation]
\label{thm:risk-aggregation}
Suppose \(\mathcal D_f\preceq\mathcal D_c\) through
\(q:V_f\to V_c\), and suppose the aggregation is information faithful. Then
for every coarse coalition \(S\subseteq V_c\),
\[
\boxed{
\Risk_{\mathcal D_c}(S)
=
\Risk_{\mathcal D_f}\bigl(q^{-1}(S)\bigr).
}
\]
Consequently,
\[
k_\varepsilon(\mathcal D_c)
\leq
k_\varepsilon(\mathcal D_f),
\]
where coalition size is counted in displayed nodes.
\end{theorem}

\begin{proof}
Information faithfulness gives
\[
\mathscr F_S^{\mathcal D_c}
=
\bigvee_{B\in S}
\bigvee_{i\in q^{-1}(B)}
\mathscr F_i^{\mathcal D_f}
=
\mathscr F_{q^{-1}(S)}^{\mathcal D_f}.
\]
The two Bayes risks optimize over the same measurable decision rules and are
therefore equal. If a fine coalition of size \(k\) is sufficient within
\(\varepsilon\), the set of coarse blocks meeting that coalition has at most
\(k\) elements and contains at least the same information. This proves the
coalition-size inequality.
\end{proof}

The theorem makes the source of a common error explicit. A minimal coalition
size of two at one resolution can become one after those two nodes are merged.
Node count is therefore a decomposition-relative statistic.

A more stable quantity measures primitive carrier cost. Let
\(w_\omega>0\) be a declared cost for each \(\omega\in\Omega\), and assign a
block the cost
\[
c(B)=\sum_{\omega\in B}w_\omega.
\]
Define the sufficient support cost
\[
\boxed{
s_\varepsilon(\mathcal D)
=
\min_{S\subseteq V_{\mathcal D}}
\left\{
\sum_{B\in S}c(B):
\Risk_{\mathcal D}(S)
\leq
\Risk_{\mathcal D}(V_{\mathcal D})+\varepsilon
\right\}.
}
\]
If \(\mathcal D_f\preceq\mathcal D_c\) is information faithful, then
\[
\boxed{
s_\varepsilon(\mathcal D_f)
\leq
s_\varepsilon(\mathcal D_c).}
\]
The fine decomposition permits every union of coarse blocks and may also
permit cheaper sub-block coalitions. Equality holds whenever an optimal fine
coalition is a union of complete coarse fibers. This weighted support measure
does not remove the need to choose primitive carriers, but it prevents a
large macro-component from being counted as a single inexpensive unit merely
because it was assigned one node label.

\topic{57}{Trivialization by Arbitrary Redescription}

The dependence on decomposition is unavoidable. Two elementary constructions
show why a certificate cannot be required to survive every mathematically
possible partition.

\begin{proposition}[Coarsening and secret-sharing trivializations]
\label{prop:decomposition-trivialization}
The following statements hold.
\begin{enumerate}
\item Any finite realization can be coarsened to a single displayed component.
 Relative to that decomposition, the full-system state is local and
 \(k_0\leq1\).
\item Let \(Z\) be uniform on a finite abelian group \(\mathsf G\). Let
 \(K\) be an independent uniform group element and define
 \[
 X_1=K,
 \qquad
 X_2=Z-K.
 \]
 Then
 \[
 I(Z;X_1)=I(Z;X_2)=0,
 \qquad
 H(Z\mid X_1,X_2)=0.
 \]
 Under logarithmic loss,
 \[
 \Risk(\{1\})=\Risk(\{2\})=\log|\mathsf G|,
 \qquad
 \Risk(\{1,2\})=0,
 \]
 so the declared two-node representation has \(k_0=2\) and joint advantage
 \(\log|\mathsf G|\), even when a single centralized encoder generated both
 shares.
\end{enumerate}
\end{proposition}

\begin{proof}
The first statement follows by taking the indiscrete partition
\(\Pi=\{\Omega\}\). For the second, conditioning on any value of \(Z\) leaves
\(K\), and hence \(X_1\), uniform. Translation invariance of the uniform law
also leaves \(X_2=Z-K\) uniform, so each share is independent of \(Z\). The
pair reconstructs \(Z=X_1+X_2\), which gives the entropy and Bayes-risk
identities.
\end{proof}

The secret-sharing example is a valid distributed code. What it does not show
is that the process producing or using the code is causally distributed across
the declared carriers. If one hidden controller creates both shares, stores
the key, or jointly selects their updates, that controller belongs in the
causal inventory. Coalition irreducibility and matched predictive nulls remain
necessary tests, but their interpretation is conditional on a defensible
interface model.

This distinction suggests two levels of claim:
\begin{description}
\item[Encoding-relative relationality.] No declared local view is sufficient.
\item[Carrier-relative relationality.] The declared views also correspond to
admissible, separately tracked causal interfaces with no omitted common
controller.
\end{description}
The stronger claim requires the decomposition audit developed next.

\topic{58}{Admissible Decomposition Families}

No purely formal rule can identify the uniquely correct decomposition of every
physical or computational system. The practical goal is instead to define a
family of decompositions that are defensible for the question being asked and
to report how the certificate changes across that family.

\begin{definition}[Admissible decomposition family]
\label{def:admissible-decomposition-family}
An admissible family \(\mathfrak D\) is a predeclared collection of
decompositions of the same carrier presentation satisfying the following
audit conditions.
\begin{enumerate}
\item \textbf{Carrier continuity.} Each component corresponds to a process,
 artifact, or mechanism that can be tracked through the relevant time scale.
\item \textbf{Access locality.} The local state contains only information
 available through the declared component and its recorded incoming channels.
\item \textbf{Channel completeness.} Shared memories, buses, supervisors,
 caches, and synchronized external cues are represented explicitly.
\item \textbf{Intervention separability.} At least some component mechanisms
 or channels can be altered without directly rewriting the entire system.
\item \textbf{Target independence.} The decomposition and its alternatives
 are fixed without using held-out outcomes to manufacture a favorable result.
\item \textbf{Resolution relevance.} A proposed split is retained only when
 it changes access, intervention, reconstruction, or transfer structure at the
 declared measurement resolution.
\item \textbf{Alternative coverage.} The family contains the principal
 plausible merges, splits, and artifact assignments that a skeptical analyst
 would consider.
\end{enumerate}
\end{definition}

Let \(m_q(W;\mathcal D)\) be the signed margin of gate \(q\), oriented so that
\[
m_q(W;\mathcal D)>0
\]
means that the gate passes after its uncertainty correction. Define the
decomposition envelope
\[
\underline m_q
=
\inf_{\mathcal D\in\mathfrak D}m_q(W;\mathcal D),
\qquad
\overline m_q
=
\sup_{\mathcal D\in\mathfrak D}m_q(W;\mathcal D),
\]
and its sensitivity width
\[
\Delta_q^{\mathrm{dec}}
=
\overline m_q-\underline m_q.
\]

\begin{definition}[Decomposition-robust certificate]
\label{def:decomposition-robust-certificate}
A workspace passes the \(\mathfrak D\)-robust certificate at level \(L\) when
\[
\boxed{
\operatorname{Certify}_{L,\mathfrak D}^{\mathrm{rob}}(W)
=
\bigwedge_{q\in\mathcal J(L)}
\left\{
\inf_{\mathcal D\in\mathfrak D}
 m_q(W;\mathcal D)>0
\right\}.
}
\]
A pass for one selected decomposition is a \emph{decomposition-specific}
claim. A pass throughout \(\mathfrak D\) is a \emph{family-robust} claim. A
claim is \emph{decomposition fragile} when at least one required gate changes
sign inside the admissible family.
\end{definition}

Universal invariance over all partitions is neither possible nor desirable,
as Proposition~\ref{prop:decomposition-trivialization} shows. The family
\(\mathfrak D\) must be broad enough to challenge the favored description and
narrow enough to exclude partitions with no operational meaning.

\topic{59}{Stability Under Decomposition Transport}

Exact faithful aggregation preserves the structural core, but practical
alternative decompositions usually introduce approximation. Their certificate
can still be controlled by transporting a reference workspace and comparing
gate margins.

Fix a reference decomposition \(\mathcal D_0\). For each
\(\mathcal D\in\mathfrak D\), let
\[
\Phi_{\mathcal D}:W^{\mathcal D_0}\longrightarrow W^{\mathcal D}
\]
be a declared workspace correspondence and let
\(\delta_{\mathcal D}\) measure its structural and causal identity defect.
Assume gate \(q\) has a transport bound
\[
\left|
 m_q(\Phi_{\mathcal D}W;\mathcal D)
 -m_q(W;\mathcal D_0)
\right|
\leq
L_q\delta_{\mathcal D}+b_q(\mathcal D),
\]
where \(L_q\) is a sensitivity constant and \(b_q\) records model mismatch not
captured by the identity distance.

\begin{theorem}[Decomposition-transport robustness]
\label{thm:decomposition-robustness}
If, for every required gate \(q\),
\[
\boxed{
 m_q(W;\mathcal D_0)
>
\sup_{\mathcal D\in\mathfrak D}
\left(
L_q\delta_{\mathcal D}+b_q(\mathcal D)
\right),
}
\]
then \(W\) passes the \(\mathfrak D\)-robust certificate. Conversely, if for
some required gate
\[
 m_q(W;\mathcal D_0)
+
\sup_{\mathcal D\in\mathfrak D}
\left(
L_q\delta_{\mathcal D}+b_q(\mathcal D)
\right)
<0,
\]
then that gate fails throughout the family.
\end{theorem}

\begin{proof}
For every \(\mathcal D\in\mathfrak D\), the transport bound gives
\[
 m_q(\Phi_{\mathcal D}W;\mathcal D)
\geq
m_q(W;\mathcal D_0)
-
L_q\delta_{\mathcal D}
-
b_q(\mathcal D).
\]
The first hypothesis makes the right-hand side positive uniformly in
\(\mathcal D\), which proves robust passage of every required gate. The second
statement follows from the corresponding upper bound.
\end{proof}

When \(b_q\equiv0\), a convenient certificate radius is
\[
\boxed{
r_{\mathrm{dec}}(W)
=
\min_{q:L_q>0}
\frac{m_q(W;\mathcal D_0)}{L_q}.
}
\]
Every admissible family satisfying
\[
\sup_{\mathcal D\in\mathfrak D}\delta_{\mathcal D}
<
r_{\mathrm{dec}}(W)
\]
preserves the sign of all required gate margins. In applications, the
transport terms can be bounded using earlier results:
\begin{center}
\small
\begin{tabularx}{\textwidth}{@{}>{\bfseries\raggedright\arraybackslash}p{0.22\textwidth}>{\raggedright\arraybackslash}X@{}}
\toprule
Gate family & Relevant transport control \\
\midrule
Structural coherence & Transfer-Laplacian projector drift and spectral gap. \\
Predictive alignment & Divergence between transported conditional-law states. \\
Causal efficacy & Distance between matched intervention kernels. \\
Relationality & Coalition-risk distortion and matched-null preservation error. \\
Reconstruction & Perturbation of lower frame bounds and handoff error. \\
Persistence and identity & The coordinatewise identity profile from Part~XII. \\
\bottomrule
\end{tabularx}
\end{center}

The theorem is intentionally margin based. A result barely above threshold at
one decomposition should not be described as robust when a small, defensible
change of interface makes it disappear.

\topic{60}{Decomposition Audit and Reporting}

A complete decomposition audit should proceed in the following order.
\begin{enumerate}
\item Inventory the primitive carriers at the chosen temporal and physical
 resolution.
\item Record every shared resource and distinguish local state from incoming
 information.
\item Choose a reference decomposition from independently manipulable and
 temporally trackable interfaces.
\item Construct an admissible family containing the principal plausible
 merges, splits, and artifact assignments.
\item For every alternative, specify the correspondence to the reference
 workspace and validate its manipulation checks.
\item Recompute the structural, predictive, causal, coalition, reconstruction,
 and identity margins.
\item Report the full decomposition envelope, not only the most favorable
 member of the family.
\end{enumerate}

The report should state:
\[
\begin{aligned}
&\text{reference decomposition and audit resolution},\\
&\text{admissible family and reasons for excluded partitions},\\
&\text{hidden-resource and common-controller checks},\\
&\text{faithful-aggregation maps and identity defects},\\
&\text{gatewise minima, maxima, and sensitivity widths},\\
&\text{strongest decomposition-specific and family-robust conclusions}.
\end{aligned}
\]

Four controls are especially important:
\begin{enumerate}
\item merge components connected by an uninterruptible common controller and
 test whether the relational claim collapses;
\item split a known centralized controller into complementary local views and
 verify that the protocol rejects the resulting secret-sharing appearance of
 distributedness;
\item expose or remove shared caches, external cues, and common memory that
 were absent from the favored graph; and
\item repeat the frozen certificate across the complete admissible family.
\end{enumerate}

For the motivating human--AI collaboration, this analysis changes the form of
the empirical question. It is not enough to ask whether the human, model,
executor, and artifact store form a workspace under one natural description.
The stronger question is whether the same operative state and its causal role
survive plausible alternatives, such as merging model and executor, separating
conversation state from persistent artifacts, or making shared orchestration
channels explicit. Those measurements have not yet been performed. The case
therefore remains decomposition-specific and observational.
